\documentclass[9pt,a4paper]{article}
\usepackage[left=0.55cm,right=0.55cm,top=0.25cm,bottom=0.25cm]{geometry}
\usepackage{enumitem}
\usepackage{hyperref}
\usepackage{xcolor}
\usepackage{titlesec}
\usepackage{fontspec}
\definecolor{darkblue}{HTML}{0F172A}
\definecolor{accentblue}{HTML}{2563EB}
\hypersetup{colorlinks=true,urlcolor=accentblue,linkcolor=darkblue}
\setlength{\parindent}{0pt}
\setlength{\parskip}{0pt}
\renewcommand{\familydefault}{\sfdefault}
\linespread{0.94}
\raggedbottom
\titleformat{\section}{\normalsize\bfseries\color{darkblue}}{}{0pt}{}[\titlerule]
\titlespacing*{\section}{0pt}{0pt}{0pt}
\setlist[itemize]{leftmargin=1.15em,itemsep=0pt,topsep=0pt,parsep=0pt,partopsep=0pt}
\begin{document}

\begin{center}
{\LARGE \textbf{ANDRÉ NICOLAS SILVA}}\\[-2pt]
{\normalsize \textbf{UX / Product Designer | UX Engineer | UI, Produto e IA}}\\[1pt]
\footnotesize Maceió, AL \quad • \quad \href{mailto:devandrenicolas@gmail.com}{devandrenicolas@gmail.com}\\[-1pt]
\textbf{\href{https://portfolio-andrenicolas.vercel.app}{Portfólio}} \quad • \quad \href{https://github.com/devAndreNicolas}{GitHub} \quad • \quad \href{https://www.linkedin.com/in/devandrenicolas}{LinkedIn}
\end{center}

\vspace{-3pt}
\section*{RESUMO PROFISSIONAL}
\footnotesize
Profissional de produto digital com base técnica em engenharia de software e foco em UX, UI e experiência centrada no usuário. Atua na interseção entre problema, jornada, interface, regras de negócio e implementação, transformando requisitos e necessidades de usuários em fluxos claros, componentes reutilizáveis e experiências tecnicamente viáveis.

Experiência em produtos SaaS de privacidade, compliance e auditoria digital utilizados por 10+ empresas, participando de todo o ciclo de entrega: entendimento do contexto, análise de requisitos, definição de escopo, decisões de UX e arquitetura, implementação frontend/full stack, refinamento, debugging, qualidade, produção e suporte. A combinação de UX engineering e conhecimento de produto permite colaborar de forma concreta com design, produto, engenharia e stakeholders.

Também desenvolve produtos próprios com jornadas de convite, colaboração, checkout, simulação financeira, operação offline-first e interfaces responsivas. Possui repertório em User-Centered Design, Design Thinking, Double Diamond, discovery, user research, benchmarking, desk research, user flow, user journey, wireframes, prototipagem e testes de usabilidade como métodos e conhecimentos de design de produto; as experiências profissionais documentadas estão mais fortes em UX engineering, interface, fluxos, implementação e iteração do que em pesquisa formal ou testes A/B.

Utiliza IA generativa e ferramentas AI-native para explorar soluções, compreender codebases, revisar alternativas, investigar problemas, gerar testes e iterar sobre funcionalidades, mantendo julgamento técnico, clareza para o usuário e responsabilidade sobre a implementação.

\section*{COMPETÊNCIAS EM UX, UI E PRODUTO}
\footnotesize
\textbf{Experiência do usuário:} UX; UI; User-Centered Design (UCD); design centrado no usuário; product thinking; user flows; user journeys; jornadas digitais; arquitetura da informação; clareza de conteúdo e estados; usabilidade; acessibilidade; design responsivo; mobile-first; redução de atrito; resolução de problemas; empatia com diferentes contextos de uso.\\
\textbf{Pesquisa e discovery:} User Research; entrevistas e investigação de necessidades como conhecimentos de processo; Benchmarking; Desk Research; análise de referências; definição de problemas; hipóteses de valor; Continuous Discovery como referência de trabalho; priorização de oportunidades; feedback e iteração.\\
\textbf{Métodos de design:} Design Thinking; Double Diamond; ideação; definição; exploração; prototipagem; wireframes; protótipos interativos; documentação de fluxos; revisão de design; handoff e alinhamento com engenharia; testes de usabilidade e testes A/B como métodos conhecidos, sem alegar experiência formal não documentada.\\
\textbf{Produto e colaboração:} requisitos de produto; regras de negócio; definição de escopo; stakeholders; trade-offs; design reviews; colaboração multidisciplinar; comunicação escrita e verbal; organização; pensamento crítico; criatividade; adaptabilidade; resiliência; orientação técnica durante a implementação.\\
\textbf{UI engineering e design systems:} Angular; React; Next.js; Astro; TypeScript; JavaScript; HTML; CSS; Web Components; Lit; Shadow DOM; Tailwind CSS; componentização; design systems; bibliotecas de componentes; tokens visuais; estados de interface; responsive design; acessibilidade; SEO; performance; Lighthouse.\\
\textbf{Qualidade da experiência:} testes automatizados; Vitest; code review; debugging; análise de comportamento; validação de regras; performance; produção; CI/CD; integração com APIs; SQL/PostgreSQL; monitoramento de estados e fluxos.\\
\textbf{IA aplicada a produto:} IA generativa; LLMs; Genkit; contextual prompting; prompt engineering; sanitização de entradas; respostas estruturadas; Codex; Claude; Cursor; CodeRabbit; MCPs; agents; specification-driven development.

\section*{EXPERIÊNCIA PROFISSIONAL}
\textbf{MSPA — Frontend Engineer (contribuições full stack) | foco em produto e UX} \hfill \textit{Abr 2025 — Atual}\\
\textit{Plataforma SaaS para privacidade, compliance (LGPD) e auditoria digital}
\begin{itemize}
\item Participo da construção e evolução de Compass, Trust Center e superfícies de compliance utilizadas por 10+ empresas, conectando necessidades de produto, regras de negócio, restrições técnicas e experiência do usuário.
\item Transformo requisitos e problemas de produto em fluxos, estados de interface, componentes, integrações e comportamentos implementáveis, acompanhando o trabalho desde a definição técnica até refinamento, debugging, produção e suporte.
\item Contribuo para decisões de UX e arquitetura frontend em Angular, React, Astro e TypeScript, buscando clareza de navegação, consistência de interação, responsividade, acessibilidade, manutenibilidade e evolução segura.
\item Desenvolvo interfaces e componentes reutilizáveis com Angular Signals, RxJS, Web Components, Lit e Shadow DOM, considerando separação entre comportamento, estado, apresentação e regras de negócio.
\item Contribuo para um consent engine agnóstico de framework, conectando interfaces de consentimento, configurações administráveis, eventos de auditoria e Cloudflare Workers, KV, D1 e Durable Objects. O desafio de produto é tornar regras de privacidade compreensíveis, confiáveis e acionáveis.
\item Trabalho com fluxos de produtos sensíveis a confiança, compliance e auditoria, considerando feedback, consistência de estados, mensagens claras, acessibilidade, tratamento de erros e impactos da implementação na experiência.
\item Itero sobre problemas de usabilidade, performance e SEO com debugging, Lighthouse, testes automatizados, code review, análise de comportamento e colaboração com produto/engenharia.
\item Integro funcionalidades de IA generativa voltadas ao usuário com Genkit, contexto, sanitização de entradas e respostas estruturadas, considerando clareza, limites do produto, previsibilidade e confiabilidade.
\item Atuo além da camada visual quando necessário, trabalhando com serviços Go, APIs REST/internas, SQL/PostgreSQL, processamento assíncrono, eventos, SSE, integrações e automações para entregar fluxos completos.
\item Utilizo Codex, Claude, Cursor, CodeRabbit, MCPs, agents e SDD para explorar alternativas, compreender sistemas, revisar alterações, investigar problemas, gerar testes e iterar, mantendo julgamento técnico humano sobre produto e qualidade.
\item Reorganizei o fluxo de desenvolvimento durante uma fase crítica de uma equipe de 6 pessoas, coordenando diretamente 4 integrantes para melhorar clareza de escopo, organização de sprints e previsibilidade da entrega; não exerci cargo formal de gestão.
\end{itemize}

\section*{CONTRIBUIÇÕES OPEN SOURCE}
\textbf{Stoat / Revolt Platforms LTD — Frontend Engineer voluntário} \hfill \textit{2026 — Atual}\\
Contribuição frontend/UX no cliente web open source do Stoat, identificando um problema voltado ao usuário e implementando uma correção de comportamento e comunicação de interface. A pull request \href{https://github.com/stoatchat/for-web/pull/1518}{\#1518} foi integrada pelos mantenedores.

\section*{PROJETOS DE PRODUTO, UX E INTERFACE}
\textbf{FechaRacha — Metas colaborativas, convites e contribuições} \hfill \textit{Next.js, React, TypeScript, Supabase, PostgreSQL, Prisma, Resend, Zod, Vitest}\\
Modelei e implementei um produto para criação de metas em grupo, convite de participantes e acompanhamento de contribuições declaradas, orientado a um caso de uso brasileiro de coordenação via WhatsApp + Pix. A estratégia de produto documenta usuários, problema, jornada de entrada, aquisição por link, ativação, hipóteses de valor, privacidade, não-custódia e eventos de produto.
\begin{itemize}
\item Estruturei fluxos de criação, convite, entrada, participação e acompanhamento, considerando papéis, capacidades, visibilidade, estados e mensagens de negócio.
\item Conectei autenticação, persistência, notificações e regras de autorização para que a experiência refletisse corretamente o contexto de cada participante.
\item Testes comportamentais verificam ownership no limite da consulta, permissões derivadas de papel, visibilidade por membro, reutilização de convite pendente, consumo atômico do limite de convites, deduplicação de notificações e transições de status.
\item A estratégia usa hipóteses de valor, analytics com privacidade e redução de PII; preços, anúncios, gamificação e expansão permanecem hipóteses de roadmap, não resultados comprovados.
\end{itemize}

\textbf{Quebrando Fronteiras — Plataforma social e de doações} \hfill \textit{Next.js, TypeScript}\\
Refatorei fluxos críticos de usuário e comportamento mobile, identificando problemas de UX e implementando melhorias com foco em clareza, redução de atrito e conversão. O trabalho evidencia iteração de experiência orientada por contexto e analytics; não alego testes A/B específicos ou métricas de conversão não registradas.

\textbf{MSPA Trust Center / Consent Engine — Experiência de consentimento} \hfill \textit{Lit, Web Components, Shadow DOM, Cloudflare}\\
Contribuí para uma experiência de consentimento framework-agnostic, com UI, configuração administrável, enforcement, eventos de auditoria e integração com serviços edge. O trabalho exige traduzir regras de privacidade em interações compreensíveis, estados previsíveis e componentes reutilizáveis.

\textbf{Saúde em Campo — Operação de saúde offline-first} \hfill \textit{React Native, Expo, TypeScript, NativeWind, SQLite, Drizzle, TanStack Query}\\
Desenhei a arquitetura de uma experiência móvel para registros de campo sem conexão estável, considerando navegação por papel, isolamento de dados por usuário, fila de sincronização, idempotência, retries/backoff, RBAC e priorização de risco. É evidência de arquitetura e design de sistema, não de implantação clínica.

\textbf{Diário de Campo Escoteiro — PWA para registros em campo} \hfill \textit{React, TypeScript, Vite, Tailwind, PouchDB/CouchDB}\\
Desenvolvi uma aplicação instalável para registros de campo com dados locais, service worker e sincronização automática após reconexão. O projeto demonstra atenção a continuidade de uso, estados offline/online, feedback de sincronização e redução de dependência de conectividade.

\textbf{Terto Beats — Marketplace de ativos digitais} \hfill \textit{Next.js, TypeScript, Supabase, Vercel}\\
Construí uma experiência de comércio digital ponta a ponta, conectando apresentação de produtos, autenticação, estado de pedido, checkout, persistência, integrações e automações pós-compra. O projeto demonstra desenho de jornada comercial e implementação de estados críticos, sem alegar receita ou volume de clientes.

\textbf{RendeCerto — Simulador de cenários financeiros} \hfill \textit{Next.js, TypeScript, Vitest}\\
Desenvolvi uma interface para transformar cenários financeiros em regras determinísticas, comparações e resultados legíveis, com testes automatizados para cenários principais. O projeto demonstra comunicação de informação, clareza de saída e validação de comportamento, sem aconselhamento financeiro.

\textbf{Sites de comércio — Superfícies web recorrentes} \hfill \textit{Next.js 16, React 19, TypeScript, Tailwind, Vercel Analytics}\\
Entreguei três superfícies web orientadas a comércio, com deploy na Vercel, componentes reutilizáveis, responsividade e analytics. A evidência sustenta entrega técnica recorrente, não relação comercial, receita ou conversão mensurada.

\section*{FORMAÇÃO}
\textbf{Tecnólogo em Sistemas para Internet} — UNCISAL \hfill \textit{Concluído em 2026}\\
Formação superior tecnológica brasileira com base em desenvolvimento web, sistemas, programação, banco de dados e construção de produtos digitais. Não é graduação específica em Design de Produto ou Design Gráfico.

\section*{IDIOMAS E DISPONIBILIDADE}
\textbf{Português:} Nativo \quad • \quad \textbf{Inglês:} Intermediário/avançado, em aprimoramento.\\
Disponibilidade para atuação presencial em Maceió, AL, e colaboração próxima a produto, engenharia, usuários e stakeholders.

\end{document}
